\cotefile{100}
\scanpage{41}
\uline{Remarque} \deleted{\unsure{Supp.|Supposons}} Soit \deleted{\unsure{B}}$S$ \inserted{$=\unsure{\mathrm{Spm}|\mathrm{Spf}|\mathrm{Spec}}\,R$} un schéma local noeth. complet,
$\mathfrak X$ un schéma formel propre \deleted{et} plat sur $S$,
tel que $H^0(X_0,\mathcal O_{X_0})\simeq k$ (pour simplifier). Supposons
aussi \deleted{\illegible{1}} que \deleted{$\underline{\mathrm{Pic}}^0_{X_0/k}$ soit lisse, soit que}
les $\underline{\mathrm{Pic}}^0_{X_n/S_n}$ soient plats \deleted{et que} \inserted{(p. ex. $\underline{\mathrm{Pic}}^0_{X_0/k}$ lisse)} \unsure{de sorte|si bien} qu'on
peut former $\underline{NS}_{X_n/S_n}=\underline{\mathrm{Pic}}_{X_n/\unsure{k_n|S_n}}/\underline{\mathrm{Pic}}^0_{X_n/S_n}$, et
le sous-schéma $P_n=\underline{NS}_{X_n/S_n}^{+}$ (partie ample).
Les $P_n$ sont \deleted{\illegible{1}} \unsure{nets|étales|ouverts} sur les $S_n$, et par \unsure{là|la}
\unsure{compatibles|suite ils} définissent un schéma formel $\widehat P$ sur $\widehat S$,
donc que $\widehat P=\coprod_{\alpha}\widehat P_\alpha$
les $P_\alpha$ étant finis \inserted{nets} sur $S$, et connexes. Procédant
comme plus haut, on trouve que les conditions suivantes
sont équivalentes

(i) $\exists\ S'\to S$ \deleted{\unsure{propre}} \inserted{propre} surjectif, \inserted{\deleted{\illegible{1}}} avec $\widehat X_{\widehat S'}$ algébrique
et \uline{projectif} sur $S'$.

(ii) $\forall$ sous-schéma fermé \inserted{$T$ de $S$} de $\dim\leq 1$, si $T'$ est son normalisé,
alors $\widehat X_{\widehat T'}$ est algébrique et projectif sur $\widehat T'$.

(iii) $\exists$ morphisme \inserted{fini} surjectif net $S'\to S$, tel que
$\widehat X_{\widehat S'}$ soit algébrique et projectif.
% slanted
\margin{\uline{NB} Cette condition n'est bonne telle quelle que si $\underline{\mathrm{Pic}}^0_{\unsure{X_0/S_0|X_0/k}}$ est lisse; sinon \unsure{on sera} \deleted{\unsure{obligé}} obligé de prendre \inserted{(iii) et (iv)} $S'\to S$ \uline{fini et plat} (\unsure{car une polarisation|au sens f. loc. libre} \unsure{ne définit pas tjrs} un faisceau inversible\ldots}

(iv) Il existe un morphisme fini étale surjectif $S'\to S$, tel que pour tte composante irréductible
$S'_i$ de $S'$, avec str. réduite induite, on a $\widehat X_{\widehat S'_i}$ est
algébrique et projectif

\deleted{Ces conditions \unsure{impliquent la suivante}\unsure{.} \unsure{cette dernière}} De
plus, si $S'\to S$ est comme dans (iv), $\widehat X_{\widehat S'}$ est algébrique.
\deleted{\unsure{(iv$'$)} \unsure{(iv$'$)|(vi$'$)|(iv)} $\exists$ morphisme fini étale surjectif}
% slanted
\margin{\deleted{\unsure{Le}} \unsure{Les} deux dernières assertions}
% archive sheet number: 40
\scanpage{42}
EGA IV 15.6. Dans 15.6.1. (p.236 l. -18) supprimer \textquotedbl localement
noethérien\textquotedbl\ , cette hypothèse ne sert pas dans la démonstration; en revanche, dans 15.6.3. \inserted{et 15.6.4.} il faut al\deleted{\unsure{or|ro}}ors préciser qu'on suppose $f$ localement
de présentation finie, et dans \fbox{15.6.4. que de plus $f$ est quasi-séparé.}
% Hand-drawn frame, and oblique stroke between the first and third lines.

\handwritten{$\Bigg($}On notera que sous ces conditions, \deleted{la}le \deleted{proposition} corollaire 15.6.5.
devient inutile, étant déjà contenu dans ce qui précède.

IV 15.6.7. p.241, l. 2, après \textquotedbl no\deleted{\unsure{t}}ethérien\textquotedbl\ ajouter \textquotedbl ou si la condition
\handwritten{$\unsure{\mathrm{D}|1}$} de 15.6.6. est vérifiée\textquotedbl\ .

Il convient également d'ajouter quelques points d'\uline{autocritique}
(peut-être bloqués ensemble, en C) après les \textquotedbl modifications de textes),
savoir les points suivants:

1\textordmasculine) Suivant une idée de Raynaud, il apparait maintenant qu'on peut
traiter de façon nettement plus satisfaisante le co\deleted{\unsure{m}}ntenu des par.14 et
15 en utilisant l'outil de la hensélisation developpé au par.18 .
De cette façon, on arrive à des formulations et des démonstrations plus
simples, tous les résultats principaux étant de plus indépendants de toute hypothèse noethérienne. De plus, on arrive à établir aussi par cette
voie que la notion de morphisme universellement ouvert \deleted{\unsure{se compor}} de
présentation finie se comporte bien par limites projectives de schémas,
suivant le point de vue du par.8 . Nous espérons présenter une version
entièrement refondue de ces résultats dans la réédition, où \deleted{ils apparai\unsure{r|x}}
la substance des actuels par.14 et 15 se trouvera rejetée après le par.
18 (ce qui ne présente aucun inconvénient, les par. 16 et suivants
étant visiblement indépendants des actuels par.14 et 15)\unsure{/.|.}

2\textordmasculine) Nous nous répentons sincèrement de l'usage immodéré fait du
terme \textquotedbl non \uline{ramifié\textquotedbl\ au lieu de \textquotedbl net\textquotedbl} . Nous pourrions ajouter quelques
% archive sheet number: 41
\scanpage{43}
\uline{NB} Si $X_0$ est un schéma abélien \deleted{et} (ou seulement
\margin{?}\unsure{paraabélien|parabélien|pro-abélien}), alors ces conditions \deleted{impliquent} \inserted{équivalent à celle (0)} que
$\unsure{\mathfrak X|\mathcal X|\widehat X}$ soit algébrique, et dans l'énoncé des conditions
(i) à (iv) [dont alors seules (i) et (ii) sont intéressantes]
on peut supprimer la condition de projectivité.

\uline{Remarque.} Soit $S'\to S$ un \inserted{morphisme propre qui est \unsure{un}} \unsure{épimorphisme|(épimorphisme)}
effectif ($S$ loc noeth). C'est donc un morphisme
de descente effective pour les schémas relatifs plats.
Soit $X'$ plat et propre sur $S'$, muni d'une donnée de
descente à $S$, on \deleted{\unsure{veut|on}} se demande si donnée
de descente est effective.

\deleted{a) Condition locale sur $S$}

a) On peut se réduire aux localisés \inserted{$\mathcal O_{S,s}$} de $S$.

b) Moyennant une question de descente \deleted{\unsure{fi|fid}} fid.
plate, on est ramené au pb sur $\widehat{\mathcal O}_{S,s}$. Il
peut être intéressant d'avoir \deleted{des \unsure{pts}} assez de p\textsuperscript{ts}
de $X_s$, et pour ceci de faire encore une ext. finie
\deleted{\unsure{étale}} \inserted{\deleted{\unsure{étale}} $R$} de $\widehat{\mathcal O}_{S,s}$. On peut alors supposer que
\deleted{$X_s$} sur le corps résiduel $k$ de $R$, on peut
descendre, donc qu'on a un $X_0$.

% slanted
\margin{\uline{Conditions \unsure{peut|pour}} \unsure{être} \uline{\unsure{inutiles|constantes}}}
c) Supposons \inserted{la réduction précédente faite, et} que les fibres de $X'/S'$ soient
sans aut. infinitésimaux, ou que $X'$ \deleted{soit lisse sur $S'$}
ait \deleted{\illegible{1}} ``assez'' de sections \deleted{\unsure{de plus}} \deleted{\unsure{lisses sur lesquelles}}
stables par descente, le long desquelles $X'/S'$ soit
lisse. Alors \deleted{\unsure{en|par}} un argument de modules
formels permet de construire un schéma formel
$\widehat X$ sur $\widehat S$ qui descend $\widehat X'$.
% archive sheet number: 42
\scanpage{44}
\unsure{X|x}xxxxxxxxxxxxxxxxxxxx

% Long oblique cancellation across the sheet; shorter strokes through individual entries.
\deleted{est réduit et \inserted{que les corps composants de} son anneau total des fractions \deleted{est} sont des extensions
séparables de $K$. Lorsque $A$ est un anneau de valuation discrète, donc $\widehat A$
également, ces conditions signifient donc aussi que le corps des fractions $\widehat K$ de $\widehat A$ est une extension séparable de $K$, condition automatiquement
satisfaite en particulier si $K$ est de caractéristique nulle.}

% The bibliographical line is encircled, with short diagonal slashes.
\deleted{IV p.224, dans la Bibliographie 36 remplacer Amer. J. par Annals.}

\deleted{IV 8.10.3. supprimer la co\unsure{n|m}dition 4\textordmasculine) , qui est inutile.}

% Hand-drawn outline encloses both lines.
\deleted{IV 8.12.8. lire: Soient $Y$ un schéma quasi-compact et quasi-séparé,
$f\!:\!X\quad Y$ ....}

\margin{\handwritten{10.9.11.1}}
\deleted{IV, p.104, \unsure{l.|1.} -10 , au lieu de (17.9.7) lire (17.9.6).}

% The first sentence is enclosed in pencil. Blank relation positions are left open.
\deleted{IV 12.1.1. (iv) ajouter qu'on y suppose $k\quad 1$. Ajouter une condition
(vi bis) L(ensemble des $x\quad X$ tels que $\mathrm{prof}^{\cdot}_{x}(\underline F_{f(x)})\quad k$ .}

\deleted{\unsure{111111111|lllllllll} Dém. La \deleted{condition} question étant locale, on peut supposer $X$ et $Y$
affines, donc $X$ est un sous-schéma fermé d'un $Y\ t_1,\ldots,t_n$ , et on est
ramené au cas où $X=Y\ t_1,\ldots,t_n$ . Utilisant alors (10.8.3) et ( \qquad ),
on trouve que $\mathrm{prof}^{\cdot}_{x}(\underline F_{f(x)})\quad k$ équivaut à $\underline{\mathrm{Ext}}^{i}(\underline F_{f(x)},\underline O_{X_{f(x)}})_x=0_x$
pour $i\quad n-k$, et on conclut grâce à (12.3.4) (qui, comme le lecteur le
constatera sans surprise, n'utilise pas dans sa démonstration le résultat sur la profondeur rectifié qu'on se propose d'établir).}

\deleted{IV 12.1.4. \unsure{Ajouter un}}

\deleted{\uline{Corollaire} 12.1.4.1. Soient $f\!:\!X\quad Y$ et $\underline F$ comme dans 12.1.1, et soit}

\deleted{IV \unsure{14.2.13|14.4.13} p. \unsure{212|202} l.11 14.4.8. a\textquotedbl) L'équivalence de cette propriété}

\deleted{IV 14.4.8. a\textquotedbl)b)Pour formuler ce\deleted{\unsure{tt}}s hypothèses,il faut supposer que
\deleted{avec les autres}
l'ensemble des composantes irréductibles de $Y$ est localement fini.}

\deleted{IV 14.4.10. (i), p.215 l.15 remplacer \textquotedbl quelconque\textquotedbl\ par \textquotedbl dont l'ensemble
des composantes irréductibles est localement fini\textquotedbl, et l.16 remplacer
\textquotedbl type fini\textquotedbl\ par \textquotedbl présentation finie\textquotedbl.}
% archive sheet number: 43
\scanpage{45}
d) Si on suppose \inserted{dans le Pb initial} qu'il existe $S''\to S'$ \deleted{\unsure{de pr.}} \deleted{propre} \inserted{\unsure{projectif|quasi-fini}} surjectif
avec $\unsure{\widehat X_{S'}|X_{S''}|\overline X_{S'}}$ \deleted{\unsure{p}} projectif sur $S'$, alors on trouve que
pour tt $s\in S$, \deleted{\unsure{et}} tel que pour les $s'\in S'_s$,
les $\mathrm{Pic}^0_{X'_n/S'_n}$ ($S'_n=\mathrm{Sp}\ \mathcal O_{S',s'}/\mathfrak m_{s'}^{n+1}$) sont plats,
et que les $X'_{s'}$ ($s'\in S'_s$) soient sans autom. infinitésimaux,
il existe une extension finie étale $T$ de $\mathrm{Sp}\ \widehat{\mathcal O}_{S,s}$
telle que la donnée de descente \deleted{pour} $T'/T$ sur $X'_{T'}$
soit effective, et définisse donc un $X_T$ sur $T$;
(on peut \unsure{$=$|même} supposer que les \deleted{$\unsure{\mathrm{Pic}_{X}|X|\mathrm{Pic}}$} $X_{T_i}$, où les
$T_i$ sont les comp. irréductibles réduites de $T$, sont
\uline{projectifs sur} /$T_i$).

\uline{Remarque} Soit tjrs $S'\to S$ un morphisme propre
qui est un épim effectif, avec $S$ loc. noeth. Supposons
que l'on ait sur $S'$ un schéma \inserted{propre et} plat $X'/S'$,
avec un faisceau inversible $\underline L'$ sur $X'$, et une
donnée de descente sur $(X',\underline L')$ rel. à $S'\to S$.
\deleted{Alors} Supposons de plus que pour tt $s'\in S'$,
$(X'_{s'},\underline L'_{s'})$ n'ait pas d'autom. infinitésimal non nul.
Alors la donnée de descente considérée est
effective. Argument par modules formels\ldots
% archive sheet number: 44
\scanpage{46}
EGA IV 4\textordmasculine\quad Errata et Compléments.

\deleted{$\unsure{0|O}_{\mathrm{III}}$ 10.2.1. p.18 l. -6, au lieu de \textquotedbl en outre si \textquotedbl\ lire \textquotedbl si en outre\textquotedbl}

\deleted{III 4.4.12. , p.138 l.10 , au lieu de Chap V lire Chap IV 8.11.2.}

\deleted{III 3.4. p.119, supprimer les lignes 2 et 3, qui contiennent une promesse
qui ne sera pas tenue.}

% Bracket and surrounding line by hand.
\margin{\handwritten{?}}
III 4.3.6. , p. 133 l. 7 , supprimer \deleted{la relation} \textquotedbl; on a donc $K'=K\unsure{\otimes|\mathrm O}_{A}A'$ \textquotedbl,
qui est manifestement faux.

\deleted{$0_{\mathrm{IV}}$ 22.1.1. p.183 l. 12, remplacer (20.6.20) par (20.6.4) .}

\deleted{IV Sommaire p.223, l. 16 au lieu de \textquotedbl de présentation finie\textquotedbl\ lire
\textquotedbl localement de présentation finie\textquotedbl\ ; ligne 19, après \textquotedbl comme\textquotedbl\ ajouter
\textquotedbl étant\textquotedbl.}

\deleted{IV 3.1.10. \deleted{Ajouter à la proposition une partie (ii), où on} Supposer
\inserted{d'abord} seulement $f$ quasi-compact et quasi-séparé, et \deleted{où on} conclu\deleted{\unsure{t}}re que
$\mathrm{Ass}\,f_{\handwritten{*}}(\underline F)\handwritten{\subset} f(X)$, et même $\mathrm{Ass}\, f_{\handwritten{*}}(\underline F)\handwritten{\subset} f_{\handwritten{\unsure{*|\leftarrow}}}(\mathrm{Ass}\,\underline F)$ lorsque $X$ est localement noethérien; enfin, si de plus $f$ est entier, on a l'égalité dans
cette dernière inclusion.}

\deleted{IV 3 , à la page 43, ajouter après 3.2.8. un\deleted{\unsure{e}}}

\deleted{\uline{Corollaire} 3.2.9. Soit $\underline F$ un Module cohérent sur le préschéma localement
noethérien $X$. Pour toute partie \deleted{discrète} \inserted{discrète} $S$ de $X$, \inserted{i.e. loc. finie et} telle que deux points
distincts de $S$ ne so\deleted{\unsure{nt}}ient pas spécialisation l'un de l'au\deleted{\unsure{d}}tre, et \inserted{pour} tout sys\unsure{t}tème $(\underline G(s))_{s\;S}$ de Modules quotients de lo\deleted{\unsure{gg}}ngueur finie de $\underline F_s$ , considérons le Module quotient $\underline G=\underline F/\underline H$ de $\underline F$, où $\underline H$ est le noyau de l'homomorphisme \deleted{naturel} composé $\underline F\quad i_s\ (\underline F_s)\quad i_s\ (G(s))$ , $i_s\!:\operatorname{Spec} k(s)\quad S$
étant le monomorphisme canonique. Alors $\underline G$ est un Module quotient cohérent
de $\underline F$ sans cycles associés immergés. L'application qui, à tout système
$(S,(G(s))_{s\;S})$ comme ci-dessus, associé $\underline G$ , est une bijection de l'ensemble de ces systèmes sur l'ensemble des Modules quotients cohérents de}
% archive sheet number: 45
\scanpage{48}
\uline{Modèles de Néron et groupes p-divisibles.}

Dans ce qui suit $S$ est un trait de point générique $\handwritten{\eta}$ et de point
fermé $s$. Un $S$-groupe p-divisible $(G_n,\ n\handwritten{\in}N)$ sera aussi désigné par
le symbole $\underset{\handwritten{\scriptscriptstyle\sim\!\sim}}{G}$ . On \deleted{sé}\inserted{\handwritten{dés}}igne par $p$ un nombre premier fixé.

\uline{Théorème} : Soient $G$ un $S$-schéma en groupes, $S$-néronien,\handwritten{[}dont la fibre générique est extension d'une variété abélienne par un tore\handwritten{]} et $\underset{\handwritten{\scriptscriptstyle\sim\!\sim}}{H}$ un $S$-groupe
% slanted
\margin{\handwritten{[ ] inutile}}
p-divisible. Alors tout morphisme $u_{\handwritten{\eta}}:(\underset{\handwritten{\scriptscriptstyle\sim\!\sim}}{H})_{\handwritten{\eta}}\handwritten{\to} G_{\handwritten{\eta}}$ se prolonge de manière unique en un $S$-morphisme $u:\underset{\handwritten{\scriptscriptstyle\sim\!\sim}}{H}\handwritten{\to} G$ .

\uline{Lemme} 1: Supposons $S$ hensélien. Soient $G$ un $S$-schéma en groupes , $S$-néronien \deleted{et}\handwritten{,} $M$ un $S$-schéma en groupes fini et plat , $\overset{\handwritten{\sim}}{M}$ sa partie connexe,
\deleted{$u_K$} $u_{\handwritten{\eta}}:M_{\handwritten{\eta}}\handwritten{\to} G_{\handwritten{\eta}}$ un homomorphisme. Alors , pour que $u_{\handwritten{\eta}}$ se prolonge en
un $S$-morphisme $u:M\handwritten{\to} G$, il faut et il suffit que $u_{\handwritten{\eta}}\handwritten{\mid}\overset{\handwritten{\sim}}{M}_{\handwritten{\eta}}$ se prolonge
en un $S$-morphisme $\overset{\handwritten{\sim}}{u}:\overset{\handwritten{\sim}}{M}\handwritten{\to} G$ .

% slanted
\margin{\handwritten{$G$ commutatif?}}
La nécessité est clai\deleted{\unsure{t}}\unsure{r|t}e. Prouvons la suffisance. Considérons le groupe
produit $M\times_S G$ et le morphisme
\[
\begin{aligned}
v_{\handwritten{\eta}}:M_{\handwritten{\eta}}\handwritten{\times}_{\handwritten{\eta}}G_{\handwritten{\eta}}&\handwritten{\to} G_{\handwritten{\eta}}\\
(m,g)&\handwritten{\mapsto} u_{\handwritten{\eta}}(m)+g
\end{aligned}
\]
Il nous suffit de montrer que $v_{\handwritten{\eta}}$ se prolonge en un morphisme $v:M\handwritten{\times}_S G\handwritten{\to} G$,
car la restriction de $v$ au premier facteur $M$ fournira le prolongement
de $u$ . Soit $P_{\handwritten{\eta}}$ le noyau de $v_{\handwritten{\eta}}$ et $P$ l'adhérence schématique de $P_{\handwritten{\eta}}$ dans
$M\handwritten{\times}_S G$. Le lemme 1 résulte de la propriété Néronienne de $G$ et du lemme suivant:

\uline{Lemme 2}: Le $S$-schéma en groupes $M\handwritten{\times}_S G/P$ \inserted{\handwritten{existe et}} est lisse sur $S$.

Vu l'hypothèse faite sur $u_{\handwritten{\eta}}$, le morphisme $\overset{\handwritten{\sim}}{v}_{\handwritten{\eta}}=v_{\handwritten{\eta}}\handwritten{\mid}\overset{\handwritten{\sim}}{M}_{\handwritten{\eta}}\handwritten{\times}_{\handwritten{\eta}}G_{\handwritten{\eta}}$ se prolonge en un $S$-morphisme
\[
\begin{aligned}
\overset{\handwritten{\sim}}{v}:\overset{\handwritten{\sim}}{M}\handwritten{\times}_S G&\handwritten{\to} G\\
(\overset{\handwritten{\sim}}{m},g)&\handwritten{\mapsto} \overset{\handwritten{\sim}}{u}(m)+g
\end{aligned}
\]
de sorte que l'on a la suite exacte :
\[
0\handwritten{\to}\operatorname{Ker}\overset{\handwritten{\sim}}{v}\handwritten{\to}\overset{\handwritten{\sim}}{M}\handwritten{\times}_S G\handwritten{\to} G\handwritten{\to}0 \tag*{\makebox[0pt][r]{\makebox[\displaywidth][l]{(1)}}}
\]
\margin{\handwritten{$\checkmark$}}
% archive sheet number: 47
\scanpage{49}
\mspage{- 2 -}
Il en résulte que $\operatorname{Ker}(\overset{\handwritten{\sim}}{v})$ est plat \margin{\handwritten{en fait il est \unsure{isom :|isom. à} $\widetilde M$ !}} et que $\overset{\handwritten{\sim}}{M}\handwritten{\times}_S G/\operatorname{Ker}(\overset{\handwritten{\sim}}{v})$ est lisse. Par
ailleurs, vu le choix de $\overset{\handwritten{\sim}}{M}$ , Le groupe $M\times_S G/\overset{\handwritten{\sim}}{M}\handwritten{\times}_S G=M/\overset{\handwritten{\sim}}{M}$ est étale, donc
$M\times_S G/\operatorname{Ker}(\overset{\handwritten{\sim}}{v})$ est lisse , comme extension de groupes lisses. On a évidemment
$\operatorname{Ker}(\overset{\handwritten{\sim}}{v})\handwritten{\subset} P$ , de sorte que $M\times_S G/P$ est \deleted{bien} lisse , comme quotient d'un groupe
lisse.

Démonstration du théorème. L'unicité du prolongement est claire. Comme la
propriété néronienne de $G$ est conservée par passage à l'hensélisé, on peut
supposer $S$ hensélien. Quitte alors à remplacer $\underset{\handwritten{\scriptscriptstyle\sim\!\sim}}{H}$ par sa partie connexe, on
peut , compte tenu du lemme 1 , supposer \deleted{\unsure{$H$}} $\underset{\handwritten{\scriptscriptstyle\sim\!\sim}}{H}$ connexe.

a) Cas où $G$ est semi-stable. Soit $\underset{\handwritten{\scriptscriptstyle\sim\!\sim}}{G}_{\handwritten{\eta}}$ le groupe p-divisible défini par $G_{\handwritten{\eta}}$
et soit $\underset{\handwritten{\scriptscriptstyle\sim\!\sim}}{G}'$ le $S$-groupe p-divisible obtenu à partir du Schéma en groupes $G^{\unsure{o|0|\circ}}$.
% slanted
\margin{\handwritten{``donc'' abusif}}
% «On sait» is circled; the short OK is on its lower left.
\margin{\handwritten{\uline{OK}}}
On sait que $\underset{\handwritten{\scriptscriptstyle\sim\!\sim}}{G}''_{\handwritten{\eta}}=\underset{\handwritten{\scriptscriptstyle\sim\!\sim}}{G}_{\handwritten{\eta}}/\underset{\handwritten{\scriptscriptstyle\sim\!\sim}}{G}'_{\handwritten{\eta}}$ est \inserted{\handwritten{un}} groupe p-divisible \deleted{non ramifié} \inserted{\handwritten{étale}} sur $S$ , \deleted{donc} \handwritten{et}
se prolonge en un $S$-groupe p-divisible , $\underset{\handwritten{\scriptscriptstyle\sim\!\sim}}{G}''$ , étale sur $S$. D'après \uuline{Tate},
le morphisme composé $v_{\handwritten{\eta}}:\underset{\handwritten{\scriptscriptstyle\sim\!\sim}}{H}_{\handwritten{\eta}}\overset{u_{\handwritten{\eta}}}{\handwritten{\longrightarrow}}\underset{\handwritten{\scriptscriptstyle\sim\!\sim}}{G}_{\handwritten{\eta}}\handwritten{\to}\underset{\handwritten{\scriptscriptstyle\sim\!\sim}}{G}''_{\handwritten{\eta}}$ se prolonge en un $S$-morphisme
$v:\underset{\handwritten{\scriptscriptstyle\sim\!\sim}}{H}\handwritten{\to}\underset{\handwritten{\scriptscriptstyle\sim\!\sim}}{G}''$. Comme $\underset{\handwritten{\scriptscriptstyle\sim\!\sim}}{H}$ est connexe et $\underset{\handwritten{\scriptscriptstyle\sim\!\sim}}{G}''$ étale , $v$ est nul. \deleted{P} Il en résulte
que $u_{\handwritten{\eta}}$ se factorise à travers $\underset{\handwritten{\scriptscriptstyle\sim\!\sim}}{G}'_{\handwritten{\eta}}$, donc toujours d'après \uuline{Tate}, se prolonge en un $S$-morphisme $u:\underset{\handwritten{\scriptscriptstyle\sim\!\sim}}{H}\handwritten{\to}\underset{\handwritten{\scriptscriptstyle\sim\!\sim}}{G}'$ c.q.f.d.

% slanted
\margin{\handwritten{? Cas où $G_{\unsure{\eta}}$ est une V.A. seulement !}}
\margin{\handwritten{OK}}
b) Cas général. On sait qu'il existe une \deleted{\unsure{s}}extension de traits $S'\handwritten{\to} S$\handwritten{,}
finie et plate , telle que , si $\handwritten{\eta}'$ est le point générique de $S'$, le $S'$
modèle de Néron $\handwritten{\mathscr G}'$ de $G_{\handwritten{\eta}'}$ soit semi\inserted{\handwritten{-}}stable . D'après a) , le morphisme
$u_{\handwritten{\eta}'}:\underset{\handwritten{\scriptscriptstyle\sim\!\sim}}{H}_{\handwritten{\eta}'}\handwritten{\to}G_{\handwritten{\eta}'}$ se prolonge en un $S'$-morphisme $u':\underset{\handwritten{\scriptscriptstyle\sim\!\sim}}{H}_{S'}\handwritten{\to}\handwritten{\mathscr G}'$ . Rappelons
que le $S$-modèle de Néron de $G_{\handwritten{\eta}}$ , c'est à dire $G$ , s'obtient à partir de $\handwritten{\mathscr G}'$
de la façon suivante (cf. note sur les modèles de Néron prop. 2.5):

On considère le $S$-groupe $S$-néronien $\handwritten{\mathscr G}=\handwritten{\prod}_{S'/S}(\handwritten{\mathscr G}'/S')$. Le groupe $G_{\handwritten{\eta}}$
s'identifie canoniquement à un sous-groupe de $\handwritten{\mathscr G}_{\handwritten{\eta}}$ ; soit $G_1$ l'adhérence
schématique dans $\handwritten{\mathscr G}$ de $G_{\handwritten{\eta}}$ . Alors $G$ est le \textquotedbl lissifié \textquotedbl\ de $G_1$ .
% archive sheet number: 48
\scanpage{50}
\mspage{- 3 -}
Ceci étant, il résulte de la définition de $\handwritten{\prod}_{S'/S}$ et de l'existence du
prolongement $u'$ , que $u_{\handwritten{\eta}}$ se prolonge en un $S$-morphisme $u_1:\underset{\handwritten{\scriptscriptstyle\sim\!\sim}}{H}\handwritten{\to}G_1$.
Il reste à voir que $u_1$ se relève en un $S$-morphisme $u:\underset{\handwritten{\scriptscriptstyle\sim\!\sim}}{H}\handwritten{\to}G$ \deleted{$g$} . Comme $\underset{\handwritten{\scriptscriptstyle\sim\!\sim}}{H}$ est connexe, il est loisible de supposer que $p$ est la caractéristique
résiduelle de $S$. Or on a le :

\uline{Lemme 3} : Si $f:G\handwritten{\to}G_1$ désigne le morphisme canonique\inserted{\handwritten{qui induit isom. sur fibres gén.}}, il existe $n\handwritten{\geq}0$,
tel que l'élévation à la puissance $p^n$ dans $G_1$ se factorise en :
\[
G_1\overset{g}{\handwritten{\longrightarrow}}G\overset{f}{\handwritten{\longrightarrow}}G_1.
\]
(ceci a été démontré au cours du séminaire de l'an dernier).

On a alors le diagramme commutatif suivant, dans le quel la première
colonne est exacte :
% Arrows are drawn by hand. Letter and number labels are typed except the outer p^n.
\diagram{\begin{tikzcd}[ampersand replacement=\&,row sep=small,column sep=large]
0\arrow[d] \& \\
H_n\arrow[d] \& \\
\underset{\handwritten{\scriptscriptstyle\sim\!\sim}}{H}\arrow[dd,"p^n"']\arrow[r,"u_1"] \& G_1\arrow[d,"g"]\arrow[dd,bend left=45,"\handwritten{p^n}"] \\
\& G\arrow[d,"f"] \\
\underset{\handwritten{\scriptscriptstyle\sim\!\sim}}{H}\arrow[d]\arrow[ur,dashed,"u"]\arrow[r,"u_1" description] \& G_1\\
\handwritten{0}\&
\end{tikzcd}}

Le morphisme $\deleted{\unsure{f}}\handwritten{g}\,u_1$ s'annule évidemment sur $H_n$ puisqu'il en est ainsi sur
la fibre générique, donc il se factorise à travers $H=H/H_n$ , d'où l'existence du relèvement $u$.

\uline{Remarques} : La démonstration utilise à deux reprises le théorème de Tate ,
qui en toute rigueur n'est démontré que lorsque $\handwritten{\eta}$ est de caractéristique
0 . La première fois , le groupe d'arrivée $G''$ est étale et dans ce cas le
théorème de Tate n'est pas très difficile. Par ailleurs , je sais maintenant
prouver le théorème de \handwritten{T}ate dans le cas où la partie connexe , sur la fibre
générique, dans le groupe d'arrivée, est de type multiplicatif.

\uline{Corollaire} .Soit $G$ un \deleted{$S$}\inserted{$\handwritten{\eta}$}-groupe \deleted{$S$-néronien,} dont la fibre générique est exten%
% archive sheet number: 49
\scanpage{51}%
\mspage{- 4 -}%
% slanted: marginal comment on the quoted words.
sion d'une variété abélienne par un tore. Alors , pour que $G$ \deleted{es} ait \handwritten{``}bonne
réduction\handwritten{''} \margin{\handwritten{expliciter ce que ça veut dire}} sur $S$\inserted{\handwritten{de car. rés. $p>0$}}, il faut et il suffit que le groupe p-divisible $\underset{\handwritten{\scriptscriptstyle\sim\!\sim}}{G}$ se
prolonge en un $S$-groupe p-divisible .
% archive sheet number: 50
\scanpage{52}
\uline{Lemme} Soient $\deleted{\unsure{k}}K$ un corps, $G$ un groupe algébrique sur $\deleted{k}K$,
$M=(M(n))$ un groupe p-divisible ($p$ un \unsure{nb|nt} premier),
$\varphi:M\unsure{\to|-}G$ un hom de groupes. \inserted{On suppose que $\mathrm{car}\,K=p$, ou \deleted{que} $K$ parfait, et que $G$ ne contient pas de ss-groupe isomorphe à $G_a$.} Alors \inserted{\deleted{a) Si $\mathrm{car}\,K\ne p$,}} il existe
un sous-schéma en groupes $H$ de $G$, extension
d'une VA par un tore, tel que $\varphi$ se factorise
par $H$. \deleted{b) Si $\mathrm{car}\, K=$}

\uline{Dém} On a $\varphi=(\varphi(n))$, avec
\[
\varphi(n):M(n)\longrightarrow G
\]
Soit $Z_n=\mathrm{Centr}\,\varphi(n)$. Les $Z_n$ sont des sous-schémas en groupes formant une suite décroissante,
donc stationnaire, soit $Z$ la valeur commune.
Évidemment $\varphi$ se factorise par \deleted{$\unsure{G}$} $Z$, et
même par $\mathrm{Centr}(Z)$. Cela nous ramène au
cas où $G$ est commutatif. Evidemment
\deleted{\unsure{on a}} l'hom composé $M\unsure{-|\to}G-G/G^0$ est nul, donc
on peut supposer $G$ connexe.

\uline{1\textsuperscript{er} Cas} $p=\mathrm{car}. K$. On a évidemment, pour tt
entier $\unsure{\nu|v}\geq 0$, que $\varphi$ se factorise par $p^\nu G$.
Or pour $\nu$ grand, on sait \deleted{$p^\nu G$} que $p^\nu G$ est
une extension d\unsure{\textsuperscript{n}|'} une VA par un tore (à dégager
en lemme, pour tt groupe algébrique commutatif
sur corps de car. $p>0$).

\uline{2\textsuperscript{ème} Cas} \deleted{$\unsure{p\ne\mathrm{car}\,K}$. Les $M(n)$ sont alors étales.} \inserted{\deleted{\unsure{car|ou}} $K$ parfait, $G$ n'a pas de ss-grpe \unsure{isom :|isom. à} $G_a$.}
Alors on sait que $G$ est ext. d'une VA. par un tore !
% archive sheet number: 51
\scanpage{53}
\deleted{Le \inserted{\deleted{\unsure{groupe algébrique}}} groupe engendré par les $\varphi(M(n))$ est donc défini,
c'est un ss-grpe algébrique lisse de $G$. Donc}

Supposons maintenant que $G$ soit un schéma en
groupes Néronien sur un trait $S$ de corps des
fonctions $K$, soit $M$ un grpe p-divisible, et
$\varphi_K:M_K\unsure{-|\to}G_K$. Je dis que $\varphi_K$ se prolonge en
$\varphi:M\to G$ (évidemment unique).

Clair si $p\ne\mathrm{car}.$ résiduelle.

Supposons donc $p=\mathrm{car}.\,k$. Alors, ou bien $\mathrm{car}.\,K=p$,
ou bien $\mathrm{car}.\,K=0$ (donc $K$ parfait); en tous cas,
$G_K$ n'a pas de ss-groupe \unsure{isom. à|isom. :} $G_a$, puisque
$G_K$ admet un modèle de Néron. Donc on
peut appliquer le lemme précédent, \deleted{Il existe} d'où
un $H_K\subset G_K$ par lequel se factorise $\varphi$.
Comme $H_K$ est ext. d'une VA par un tore,
$H_K$ admet un modèle de Néron $H$, et $H_K\overset{i_K}{\hookrightarrow}G_K$
est \unsure{induit|induite} par $H\xrightarrow{i}K$. Le th. de Raynaud
pour $H$ \unsure{implique|indique} que $\psi_K:M_K\to H_K$ se
prolonge en $\psi:M\longrightarrow H$. On gagne !
% archive sheet number: 52
